\section{Positive finite-band inversion and microscopic conditioning}
\label{sec:positive-microscopic-conditioning}
A sharp truncation of a Fourier integral need not define a positive
measure. For conditioning it is useful to retain positivity before
normalization. We construct a finite-band likelihood on the original
probability space and control it uniformly against every bounded
trajectory statistic. Only the unweighted law is regularized; the
trajectory statistic is never differentiated or smoothed.

\subsection{A positive kernel and a source-to-density estimate}
Use the real-line transform $\widehat f(b)=\int e^{ibt}f(t)\,dt$.
Define, with the continuous value at zero,
\begin{equation}\label{eq:positive-bandlimited-kernel}
 k(t)=\frac3{8\pi}\left(\frac{\sin(t/4)}{t/4}\right)^4,
 \qquad k_B(t)=B k(Bt),\qquad B>0.
\end{equation}
\begin{lemma}[Kernel normalization and moments]
\label{lem:positive-kernel-moments}
The kernel $k$ is even and nonnegative, $\int k=1$, and
\begin{equation}\label{eq:positive-kernel-transform}
 \widehat k(b)=\eta(b)=
 \begin{cases}
 1-6b^2+6|b|^3,& |b|\le1/2,\\
 2(1-|b|)^3,&1/2\le|b|\le1,\\
 0,&|b|\ge1.
 \end{cases}
\end{equation}
Moreover $\int t k(t)\,dt=0$, $\int t^2k(t)\,dt=12$, and
$\kappa_1:=\int|t|k(t)\,dt\le\sqrt{12}$.
\end{lemma}
\begin{proof}
The identity
$\sin(t/4)/(t/4)=2\int_{-1/4}^{1/4}e^{iut}\,du$
expresses the fourth power by the fourfold convolution of the interval
indicator. Evaluating that elementary convolution gives
\eqref{eq:positive-kernel-transform}, including the normalization
$\eta(0)=1$. The kernel decays as $O((1+|t|)^{-4})$ and hence has a
finite second moment. Differentiating its transform twice at zero
therefore gives $\int t^2k(t)\,dt=-\eta''(0)=12$.
Evenness gives the first signed moment; Cauchy--Schwarz gives the
absolute first-moment bound.
\end{proof}

\begin{lemma}[An all-label interpolation inequality]
\label{lem:all-label-interpolation}
Let $q_\ell\in W^{2,1}(\R)$ have compact support, and suppose
$A_j=\sum_\ell\|q_\ell^{(j)}\|_1<\infty$ for $j=0,2$.
Then
\begin{equation}\label{eq:all-label-interpolation}
 A_1\le2\sqrt{A_0A_2},\qquad
 \sum_\ell\|q_\ell\|_\infty\le A_1/2,
 \qquad
 \sum_\ell\|q_\ell-k_B*q_\ell\|_1\le6A_2/B^2.
\end{equation}
The same statements hold for real or complex functions.
\end{lemma}
\begin{proof}
Taylor's formula in $W^{2,1}$, summed over labels, gives for every $h>0$
\[
 A_1\le 2A_0/h+hA_2/2.
\]
Taking $h=2\sqrt{A_0/A_2}$ proves the first inequality when both
quantities are positive. If either vanishes the claim follows by
letting $h$ tend to zero or infinity. Each component tends to zero at
both ends of the line. Integrating its derivative from either end
bounds its supremum by half its total variation, proving the second
inequality. Finally
\[
 \|q_\ell(\cdot-y)-q_\ell+yq'_\ell\|_1
                       \le y^2\|q''_\ell\|_1/2.
\]
Integrate against $k_B(y)$, use its zero signed first moment and second
moment $12/B^2$, and sum. All sums and integrals are justified by the
nonnegative bounds just displayed.
\end{proof}

Here and below $\|\cdot\|_{\TV}$ is the total variation norm of a signed
measure, equivalently the supremum of its integrals against measurable
complex functions of modulus at most one. Thus it is the $L^1$ norm
for densities, without a factor of one half.

\begin{theorem}[Positive spectral approximation of a mixed law]
\label{thm:positive-spectral-approximation}
Let $Z=(\mathsf L,\mathsf T)$ on a probability space
$(\Omega,\mathbb P)$ take values in $\Z^3\times\R$. Suppose its law
$\mu$ admits a positive submeasure $Q\le\mu$ with compactly supported
roof densities $q_\ell\in W^{2,1}(\R)$. Write
\[
 \delta=\|\mu-Q\|_{\TV},\qquad
 A_2=\sum_\ell\|q''_\ell\|_1<\infty.
\]
Then
\begin{equation}\label{eq:positive-full-law-density}
 p_B(\ell,t)=\frac1{(2\pi)^4}
  \int_{[-\pi,\pi]^3}\int_{-B}^B
 e^{-i(u\cdot\ell+bt)}\eta(b/B)\widehat\mu(u,b)\,db\,du
\end{equation}
is a nonnegative density of total mass one. It smooths only the roof
coordinate and preserves all three lattice coordinates exactly. The
corresponding probability $\mu_B$ satisfies
\begin{equation}\label{eq:positive-full-law-TV}
 \|\mu-\mu_B\|_{\TV}\le2\delta+6A_2/B^2.
\end{equation}
If $\mu$ has density $p$, this is the global mixed $L^1$ estimate for
$p-p_B$, including the sum over all labels.
\end{theorem}
\begin{proof}
Fourier inversion for $k_B$ and torus orthogonality give
\[
 p_B(\ell,t)
   =\mathbb E\bigl[\one_{\{\mathsf L=\ell\}}
                                      k_B(t-\mathsf T)\bigr].
\]
The frequency integral is absolutely convergent, since the band is
finite and $|\widehat\mu|\le1$. This identity proves positivity,
normalization and exact preservation of the lattice marginal.
Write $\mu=Q+E$ with $E\ge0$ and $\|E\|_{\TV}=\delta$.
Roof convolution is a contraction on finite measures, so
$\|E-k_B*E\|_{\TV}\le2\delta$. The last estimate of
Lemma~\ref{lem:all-label-interpolation} bounds the remaining difference
$Q-k_B*Q$, proving \eqref{eq:positive-full-law-TV}.
\end{proof}

\subsection{A microscopic event as a finite-band likelihood}
Consider an event $A=\{Z\in\mathcal A\}$ with fibres
$\mathcal A_\ell\subset\R$, each a union of at most $J$ bounded
intervals. Empty fibres are permitted. Endpoint conventions do not
matter for the billiard laws considered below. Define on the original
probability space
\begin{equation}\label{eq:positive-event-likelihood}
 \lambda_{A,B}
       =(k_B*\one_{\mathcal A_{\mathsf L}})(\mathsf T).
\end{equation}
In particular $0\le\lambda_{A,B}\le1$. This is a weight on the same
initial trajectory, not a newly sampled section point or a change to
its physical evolution.

\begin{theorem}[Uniform weighted inversion on the original space]
\label{thm:all-selector-finite-band}
Under the hypotheses of Theorem~\ref{thm:positive-spectral-approximation},
let
\begin{equation}\label{eq:positive-event-error-budget}
 E_B(J)=\delta+2\kappa_1J\sqrt{A_2}/B.
\end{equation}
Then
\begin{equation}\label{eq:positive-likelihood-L1}
 \|\one_A-\lambda_{A,B}\|_{L^1(\mathbb P)}\le E_B(J).
\end{equation}
For every bounded measurable trajectory statistic $W$, including a
statistic depending on $n$ or $B$,
\begin{equation}\label{eq:all-selector-error}
 \left|\mathbb E[W\one_A]-\mathbb E[W\lambda_{A,B}]\right|
                               \le\|W\|_\infty E_B(J).
\end{equation}
No BV, continuity or finite-record hypothesis on $W$ is required.
For $\mathcal A=\{\ell\}\times I$, the second expectation is exactly
\begin{equation}\label{eq:all-selector-Fourier-formula}
 \frac1{(2\pi)^4}\int_{[-\pi,\pi]^3}\int_{-B}^B
 e^{-iu\cdot\ell}H_I(b)\eta(b/B)
          \mathbb E\bigl[W e^{i(u\cdot\mathsf L+b\mathsf T)}\bigr]
                                                        \,db\,du.
\end{equation}
The bound is independent of the length of $I$. Finite unions of fibres
have the corresponding finite sum of these formulas.
\end{theorem}
\begin{proof}
For a set $I$ that is a union of at most $J$ intervals,
\[
 \int_\R|\one_I(t)-\one_I(t-y)|\,dt\le2J|y|.
\]
Positivity and unit mass of $k_B$ imply pointwise
\[
 |\one_I(t)-(k_B*\one_I)(t)|
    \le\int k_B(y)|\one_I(t)-\one_I(t-y)|\,dy.
\]
On the positive discarded measure $E=\mu-Q$ the integrand is at most
one, and so costs at most $\delta$. On $Q$, bound each component by
$\|q_\ell\|_\infty$ and apply the preceding translation estimate.
Summation and Lemma~\ref{lem:all-label-interpolation} give
\[
 \int|\one_{\mathcal A}-k_B*\one_{\mathcal A}|\,dQ
   \le\frac{2\kappa_1J}{B}\sum_\ell\|q_\ell\|_\infty
   \le\frac{\kappa_1J A_1}{B}
   \le\frac{2\kappa_1J\sqrt{A_2}}B,
\]
since $A_0=\|Q\|_{\TV}\le1$. This proves
\eqref{eq:positive-likelihood-L1}. It also handles countably many
nonempty fibres by summation of nonnegative bounds.

Multiplication by $W$ proves \eqref{eq:all-selector-error}; taking the
supremum over $|W|\le1$ gives precisely the total variation norm of
$\one_A\mathbb P-\lambda_{A,B}\mathbb P$.
Fourier inversion for $k_B$, followed by integration over the bounded
interval $I$ and torus orthogonality, proves
\eqref{eq:all-selector-Fourier-formula}. Fubini applies because the
frequency band and interval are finite and $W$ is bounded.
No derivative of the weighted transform is used. In particular the
proof does not claim a global $L^1$ convolution approximation uniform
in arbitrary oscillating weights $W$; it proves the stated interval-event
operator estimate instead.
\end{proof}

\begin{corollary}[Relative conditioning and selected denominators]
\label{cor:positive-microscopic-posterior}
Let $p_A=\mathbb P(A)$ and $p_{A,B}=\mathbb E\lambda_{A,B}$.
If $p_A>E_B(J)$, the probability
$\Pi_{A,B}=p_{A,B}^{-1}\lambda_{A,B}\mathbb P$ is defined and
\begin{equation}\label{eq:positive-posterior-TV}
 \|\mathbb P(\cdot\mid A)-\Pi_{A,B}\|_{\TV}
                            \le 2E_B(J)/p_A.
\end{equation}
Alternatively, $p_{A,B}>E_B(J)$ proves the rigorous denominator bound
$p_A\ge p_{A,B}-E_B(J)>0$ and the upper bound
$2E_B(J)/(p_{A,B}-E_B(J))$ in \eqref{eq:positive-posterior-TV}.

For every measurable selector $D\subset\Omega$, the same statements
hold with $A\cap D$ and the likelihood $\one_D\lambda_{A,B}$,
using the identical error budget $E_B(J)$. More generally they hold
for a nonnegative path weight bounded by one, with its corresponding
weighted denominator. No lower bound for that denominator is assumed
without being stated.
\end{corollary}
\begin{proof}
Equation~\eqref{eq:positive-likelihood-L1} bounds both the variation
of the unnormalized measures and the difference of their masses.
For two nonnegative densities $f,g$ of masses $a,b>0$,
\[
 \|f/a-g/b\|_1\le\|f-g\|_1/a+|a-b|/a.
\]
Apply this with $f=\one_A$ and $g=\lambda_{A,B}$. Positivity of the
needed masses follows from their difference being at most $E_B(J)$.
Multiplication by $\one_D$, or by a nonnegative weight at most one,
can only decrease the unnormalized $L^1$ error. This proves every
assertion with no regularity assumption on the selector.
\end{proof}

\subsection{Uniform long-time budgets for the actual billiard laws}
Let $\mathbb P_R^{\rm sec}=\nu_R^*$ and let $\mathbb P_R^{\rm eq}$
be \eqref{eq:equilibrium-source-for-inversion}. In the first case set
$Z=J_{n,R}(y)$; in the second use the same record from the actual preceding
section origin $y$, retaining the elapsed age $s$ for all path selectors.
Denote the two output densities by $p_{n,R}^{\mathrm j}$,
$\mathrm j\in\{\mathrm{sec},\mathrm{eq}\}$.

\begin{theorem}[Microscopic positive inversion with arbitrary selectors]
\label{thm:uniform-positive-microscopic-inversion}
Fix $P>0$, $J_0<\infty$ and $\kappa\ge0$. There are choices
$L_n=\lceil\lambda n\rceil$, $\varepsilon_n=\varepsilon_0n^{-d}$
and $B_n\ge1$, uniform in $R$, for which
\begin{equation}\label{eq:positive-inversion-bandwidth}
 \log B_n\le C n\log(2+n),
\end{equation}
and the following assertions hold for both initial laws. Let
$p_{n,R,B_n}^{\mathrm j}$ be the positive finite-band density
\eqref{eq:positive-full-law-density}. Then
\begin{align}
 n^2\|p_{n,R}^{\mathrm j}-p_{n,R,B_n}^{\mathrm j}\|_{L^1(\Z^3\times\R)}
        &\le C_Pn^{-P},\label{eq:uniform-positive-density}\\
 n^2\sup_{\substack{\mathcal A:\ J\le J_0n^\kappa\\ |W|\le1}}
 \left|\mathbb E_R^{\mathrm j}[W\one_{\{J_{n,R}\in\mathcal A\}}]
       -\mathbb E_R^{\mathrm j}[W\lambda_{A,B_n}]\right|
        &\le C_Pn^{-P}.\label{eq:uniform-positive-all-selectors}
\end{align}
The supremum is over the interval-fibre events defined above and all
bounded measurable statistics on the respective original probability
space. In particular singleton lattice labels and arbitrarily shrinking
roof intervals are allowed. Their Fourier representation is
\eqref{eq:all-selector-Fourier-formula}, with $\widehat\mu^{W}$ computed
under the same initial law.
\end{theorem}
\begin{proof}
Choose $\lambda>\max\{1,a/c\}+1$ and $d>32(P+5)$.
Proposition~\ref{prop:stationary-geometric-extraction} gives, for both
laws, positive extractions with
\[
 \delta_n\le C_Pn^{-P-4},\qquad
 A_{2,n}\le A_n:=1+C\exp\{C(L_n+1)\log(C/\varepsilon_n)\}.
\]
Indeed the slower stationary geometric mass is bounded by
$C n^{1/2-d/32}$, and its high-count term is exponentially small.
Increase the common $C$ to cover both derivative budgets and all
initial values of $n\ge2$. Set
\begin{equation}\label{eq:positive-explicit-bandwidth}
 B_n=n^{P+\kappa+4}\sqrt{A_n}.
\end{equation}
This proves \eqref{eq:positive-inversion-bandwidth}. Also
\[
 2\delta_n+6A_n/B_n^2\le C_Pn^{-P-4},\qquad
 \delta_n+2\kappa_1J_0n^\kappa\sqrt{A_n}/B_n
                                    \le C_Pn^{-P-4}.
\]
Apply Theorems~\ref{thm:positive-spectral-approximation} and
\ref{thm:all-selector-finite-band}, then multiply by $n^2$.
The stronger intermediate powers absorb constants and rounding.
All bandwidths here are finite; their logarithms, not the bandwidths
themselves, have order $n\log n$.
\end{proof}

\paragraph{Location of the remaining local theorem.}
The positive likelihood gives a finite-band representation of microscopic
conditioning on the original probability space, including arbitrary
bounded path selectors and the true stationary entrance bias.
It does not evaluate that finite integral by a Gaussian. Its relative
conclusions use the denominator bounds stated in
Corollary~\ref{cor:positive-microscopic-posterior}. Likewise
\eqref{eq:uniform-positive-density} is a global $L^1$ approximation,
not an essential-supremum estimate on a prescribed central window.
The full microscopic conclusion remains governed by
\eqref{eq:geometric-common-raw-ledger}: the finite complement and the
pointwise common correction must still be estimated, followed by the
original physical-event comparison at that denominator scale.
